



% \subsection{Coordination and Scaling Benchmarks}
% \label{subsec:bench_coordination}


% To evaluate coordination, communication, and scalability in multi-agent systems, we adopt the AGENTSNET benchmark~\cite{coordinationcollaborativereasoning}, a suite of distributed decision-making tasks designed around explicit, topology-constrained message passing. AGENTSNET focuses on emergent collective behavior that arises when agents must iteratively exchange local information to reach globally coherent outcomes. At its core, the benchmark probes two fundamental coordination properties: \emph{local consistency}, where neighboring agents must resolve conflicts or negotiate decisions using only partial information, and \emph{global alignment}, where the network as a whole must converge to a valid collective state under constrained communication.

% AGENTSNET comprises five canonical distributed computing tasks that collectively span a broad range of coordination phenomena. The \emph{(1) Coloring} task captures purely local coordination by requiring agents to select labels distinct from those of their neighbors. \emph{(2) Matching} introduces bilateral negotiation, where adjacent agents must mutually select one another, testing reciprocity under decentralized control. \emph{(3) VertexCover} poses a global resource-minimization problem, requiring agents to collectively cover all edges using as few selected nodes as possible. \emph{(4) LeaderElection} enforces symmetry breaking by requiring the network to designate exactly one leader acknowledged by all agents. Finally, \emph{(5) Consensus} demands that all agents converge on a shared binary value despite potentially divergent initial preferences, making it a stringent test of long-range information propagation. Together, these tasks provide a principled and complementary basis for assessing multi-step, interdependent coordination behavior under increasing system scale.

% \paragraph{Communication topology and network structure.}
% A defining characteristic of AGENTSNET is its explicit treatment of communication topology as a first-class component of the evaluation. Rather than assuming full connectivity, agents communicate only along the edges of a predefined graph, enabling coordination performance to be studied as a function of network structure. The original benchmark employs several representative graph families, including \emph{Small-World} graphs~\cite{watts1998collective}, which combine local clustering with occasional long-range shortcuts; \emph{Scale-Free} graphs~\cite{scale_free}, which exhibit hub-dominated connectivity and rapid information diffusion; and \emph{Delaunay} graphs~\cite{delaunay}, which model spatially constrained communication common in physical and robotic multi-agent systems. These topologies allow controlled analysis of how structural properties such as path length, clustering coefficient, and node centrality influence coordination dynamics and convergence behavior.

% \paragraph{Benchmark extension and orchestration normalization.}
% Building on the original AGENTSNET benchmark, we re-engineer it into a \emph{style-oriented coordination testbed} that enables systematic, large-scale comparison of multi-agent orchestration paradigms. A central challenge in cross-framework evaluation is that different frameworks impose distinct coordination structures—graph-based workflows, role-based hierarchies, or centrally mediated interaction—that are not directly comparable under the benchmark’s native topology assumptions. To address this, we redesign the execution pipeline into a modular experimentation infrastructure that decouples task semantics from coordination style.

% Specifically, we retain the original AGENTSNET graph families as the base communication substrate and introduce explicit, deterministic \emph{topology-rewriting operators} that transform identical base graphs into alternative coordination structures. These include sequential and hierarchical topologies to emulate role-based manager–worker coordination, as well as fully connected star topologies to model centrally mediated interaction patterns. Importantly, all topology transformations preserve the original node set, task definition, and initial agent states, ensuring that observed performance differences arise solely from communication structure and orchestration strategy rather than from graph generation artifacts or task redefinition.

% \paragraph{Execution protocol and experimental controls.}
% All coordination experiments follow a standardized execution protocol. For a given task, topology, and agent count \(n\), agents execute synchronous communication rounds in which each agent observes messages from its neighbors, updates its local state according to framework-specific logic, and emits new messages along permitted communication edges. Execution terminates when either a task-specific success condition is satisfied (e.g., valid coloring, global consensus, or leader agreement) or a fixed maximum number of rounds is reached. To isolate framework-level orchestration effects, we fix the underlying LLM, prompting format, initial agent states, and stopping criteria across all frameworks, varying only the coordination topology and orchestration mechanism.

% We explicitly evaluate scalability by increasing the number of agents while holding task semantics and topology family constant. Agent counts are selected to span small to large coordination regimes, enabling analysis of how convergence behavior, communication overhead, and execution cost evolve as network size grows.

% \paragraph{Metrics and recorded signals.}
% The experimental infrastructure records a comprehensive set of metrics for each run. These include task success rate, defined as the fraction of runs that satisfy the task’s correctness criteria within the round budget; rounds to convergence, measuring the number of communication rounds required to reach a valid global state; runtime, measured as wall-clock execution time; and token cost, computed as the total number of language-model tokens consumed across all agents and rounds. In addition, the system logs structured metadata describing task configuration, topology properties, agent count, and per-agent message statistics.

% \paragraph{Reproducibility and analysis support.}
% All experimental runs are executed through a centralized experiment runner with deterministic graph loading, controlled random seeds, and unified configuration files. Results are stored in a normalized schema and automatically aggregated into summary tables suitable for scalability and efficiency analysis. A complementary visualization suite supports both static plots and interactive inspection of coordination dynamics, enabling detailed analysis of convergence patterns and failure modes as system scale increases. Collectively, these extensions transform AGENTSNET from a single-framework research prototype into a reproducible, extensible, and style-aware benchmark for studying how coordination structure, communication topology, and scaling jointly shape collective behavior in contemporary multi-agent systems.




% \vspace{-3mm}
\subsection{Coordination and Scaling Benchmarks}
\label{subsec:bench_coordination}

Coordination in multi-agent systems involves information exchange, conflict resolution, and convergence under constrained communication. Existing frameworks support agent execution and task decomposition but lack native multi-round peer-to-peer communication over arbitrary topologies. As a result, coordination is largely dictated by interaction topology rather than explicit framework mechanisms. To study topology-aware coordination and scalability, we adopt AGENTSNET~\cite{coordinationcollaborativereasoning}, which analyzes emergent collective behavior under constrained message passing.\footnote{Dataset: \url{https://huggingface.co/datasets/disco-eth/AgentsNet}; \\ Code: \url{https://github.com/floriangroetschla/AgentsNet}} AGENTSNET evaluates \emph{local consistency} and \emph{global alignment} under limited connectivity.




% \begin{figure}[t]
% \centering
% \begin{tikzpicture}

% % Define a consistent style for all axes
% \pgfplotsset{
%     common style/.style={
%         title style={font=\footnotesize\bfseries, yshift=-4pt},
%         label style={font=\scriptsize},
%         tick label style={font=\tiny},
%         ymajorgrids=true,
%         grid style={dashed, gray!30},
%         axis line style={gray!60},
%         every axis plot/.append style={
%             draw=black!70,
%             line width=0.4pt
%         }
%     }
% }

% % =========================================================
% % (a) Query Length – Top Left
% % =========================================================
% \begin{axis}[
%     common style,
%     name=qlen,
%     ybar,
%     bar width=8pt,
%     width=0.48\columnwidth,
%     height=3.2cm,
%     title={(a) Query Length},
%     xlabel={Words},
%     ylabel={Count},
%     xtick={1,2,3,4,5},
%     xticklabels={0--20,20--40,40--60,60--80,80+},
%     xticklabel style={rotate=30, anchor=east},
%     ymin=0, ymax=650,
%     enlarge x limits=0.2,
%     nodes near coords,
%     nodes near coords style={font=\tiny, gray!80, /pgf/number format/fixed}
% ]
% \addplot [pattern=north east lines, pattern color=blue!40!black!60] coordinates {
%     (1,6) (2,246) (3,458) (4,52) (5,3)
% };
% \end{axis}

% % =========================================================
% % (b) APIs per Tool – Top Right
% % =========================================================
% \begin{axis}[
%     common style,
%     name=apitools,
%     at={(qlen.east)}, anchor=west, xshift=12mm,
%     ybar,
%     bar width=8pt,
%     width=0.48\columnwidth,
%     height=3.2cm,
%     title={(b) APIs per Tool},
%     xlabel={Endpoints},
%     ylabel={\# Tools},
%     xtick={1,2,3,4,5,6},
%     xticklabels={1,2--5,6--10,11--20,21--50,50+},
%     xticklabel style={rotate=30, anchor=east},
%     ymin=0, ymax=650,
%     enlarge x limits=0.15,
%     nodes near coords,
%     nodes near coords style={font=\tiny, gray!80}
% ]
% \addplot [pattern=crosshatch, pattern color=red!40!black!60] coordinates {
%     (1,436) (2,182) (3,18) (4,13) (5,8) (6,1)
% };
% \end{axis}

% % =========================================================
% % (c) Categories – FULL WIDTH Bottom
% % =========================================================
% \begin{axis}[
%     common style,
%     at={(qlen.south west)}, anchor=north west, yshift=-14mm,
%     xbar,
%     bar width=5pt,
%     width=0.98\columnwidth,
%     height=4.0cm,
%     title={(c) Solvable Queries by Category},
%     xlabel={Query Count},
%     symbolic y coords={
%         Business,Media,Social,Sports,Entertainment,
%         Video/Images,Data,Finance,Movies,Tools
%     },
%     ytick=data,
%     xmin=0, xmax=850,
%     xmajorgrids=true,
%     ymajorgrids=false,
%     nodes near coords,
%     nodes near coords style={font=\tiny, xshift=2pt, anchor=west}
% ]
% \addplot [pattern=dots, pattern color=green!40!black!60] coordinates {
%     (96,Business) (111,Media) (152,Social) (183,Sports)
%     (213,Entertainment) (244,Video/Images) (272,Data)
%     (275,Finance) (413,Movies) (746,Tools)
% };
% \end{axis}

% \end{tikzpicture}
% \vspace{-2mm}
% \caption{Overview of StableToolBench: (a) distribution of query word counts; (b) number of API endpoints available per tool; (c) distribution of queries across top API categories.}
% \label{fig:stabletoolbench-overview}
% \vspace{-5mm}
% \end{figure}



\begin{figure}[t]
\centering
\begin{tikzpicture}

% Define a consistent style for all axes
\pgfplotsset{
    common style/.style={
        title style={font=\tiny\bfseries, yshift=-8pt},
        label style={font=\tiny},
        tick label style={font=\tiny},
        ymajorgrids=true,
        grid style={dashed, gray!30},
        axis line style={gray!60},
        % Pull y-label closer to the axis
        ylabel style={font=\tiny, xshift=5pt}, 
        xlabel style={font=\tiny, yshift=2pt},
        every axis plot/.append style={
            draw=black!70,
            line width=0.4pt
        }
    }
}

% =========================================================
% (a) Query Length - Width 27%
% =========================================================
\hspace{-4mm}
\begin{axis}[
    common style,
    name=qlen,
    ybar,
    bar width=4pt,
    width=0.39\columnwidth,
    height=3.cm,
    title={(a)},
    xlabel={Words},
    ylabel={Count},
    ylabel style={xshift=-6pt, yshift=-5pt},
    xtick={1,2,3,4,5},
    xticklabels={0-20,20-40,40-60,60-80,80+},
    xticklabel style={rotate=90, anchor=east},
    ymin=0, ymax=650,
    enlarge x limits=0.2,
    nodes near coords,
    nodes near coords style={font=\tiny, gray!80, /pgf/number format/fixed, yshift=-2pt}
]

\addplot [pattern=north east lines, pattern color=blue!40!black!60] coordinates {
    (1,6) (2,246) (3,458) (4,52) (5,3)
};
\end{axis}

\hspace{2mm}
% =========================================================
% (b) APIs per Tool - Width 27%
% =========================================================
\begin{axis}[
    common style,
    name=apitools,
    at={(qlen.east)}, anchor=west, xshift=6mm, 
    ybar,
    bar width=4pt,
    width=0.39\columnwidth,
    height=3.cm,
    title={(b)},
    xlabel={Endpoints},
    ylabel={\# Tools},
    ylabel style={xshift=-6pt, yshift=-5pt},
    xtick={1,2,3,4,5,6},
    xticklabels={1,2-5,6-10,11-20,21-50,50+},
    xticklabel style={rotate=90, anchor=east},
    ymin=0, ymax=650,
    enlarge x limits=0.15,
    nodes near coords,
    nodes near coords style={font=\tiny, gray!80}
]
\addplot [pattern=crosshatch, pattern color=red!40!black!60] coordinates {
    (1,436) (2,182) (3,18) (4,13) (5,8) (6,1)
};
\end{axis}
\hspace{5mm}
% =========================================================
% (c) Categories - Width 44%
% =========================================================

\begin{axis}[
    common style,
    name=cats,
    at={(apitools.east)}, anchor=west, xshift=10mm,
    xbar,
    bar width=3.pt,
    width=0.44\columnwidth,
    height=3.2cm,
    title={(c)},
    xlabel={Query Count},
    symbolic y coords={
        Business,Media,Social,Sports,Entertainment,
        Video/Images,Data,Finance,Movies,Tools
    },
    ytick=data,
    yticklabel style={font=\tiny},
    ylabel style={xshift=0pt}, % Reset shift for the horizontal bar chart
    xmin=0, xmax=950,
    xmajorgrids=true,
    ymajorgrids=false,
    nodes near coords,
    nodes near coords style={font=\tiny, xshift=1pt, anchor=west}
]
\addplot [pattern=dots, pattern color=green!40!black!60] coordinates {
    (96,Business) (111,Media) (152,Social) (183,Sports)
    (213,Entertainment) (244,Video/Images) (272,Data)
    (275,Finance) (413,Movies) (746,Tools)
};
\end{axis}

\end{tikzpicture}
\vspace{-5mm}
\caption{Characteristics of StableToolBench, showing the distributions of (a) query length, (b) number of APIs per tool, and (c) query counts across tool categories.}
\label{fig:stabletoolbench-overview}
\vspace{-4mm}
\end{figure}

\noindent\textbf{Benchmark Structure and Conditioning Variants.}
AGENTSNET includes five coordination tasks that capture distinct coordination behaviors: \emph{Coloring} (local conflict resolution), \emph{Matching} (bilateral agreement), \emph{VertexCover} (global resource minimization), \emph{LeaderElection} (symmetry breaking), and \emph{Consensus} (long-range information propagation)~\cite{coordinationcollaborativereasoning}. Agents interact only along edges of predefined graph families, including \emph{Small-World}~\cite{watts1998collective}, \emph{Scale-Free}~\cite{scale_free}, and \emph{Delaunay}~\cite{delaunay}. This enables systematic analysis of how clustering, path length, and connectivity heterogeneity affect coordination dynamics. To reflect interaction constraints in role-based and GABM frameworks, we introduce a graph-rewriting operator that derives sequential, hierarchical, and star-shaped interaction structures from shared base graphs. In the star configuration, interaction is mediated by a centralized environment component that serves as a communication hub for all agents.

\noindent\textbf{Evaluation Setup.}
Beyond adopting AGENTSNET, MAFBench provides a unified coordination evaluation infrastructure that standardizes topology-controlled experiments across agent systems. We implement a topology rewriting engine that transforms base communication graphs into sequential, hierarchical, and centralized structures while preserving agent sets. Dedicated runners then execute identical coordination tasks under each orchestration pattern using consistent prompts, models, and task semantics. We also introduce a visualization tool to render final agent states and communication outcomes for qualitative analysis. Finally, we automate large-scale sweeps across tasks, graph families, and agent counts, and standardize metrics to enable reproducible scalability analysis.


